\chapter{記号表}\label{ch:symbols}

文字の上のドット $\dot x$ は時間微分 $\dd x/\dd t$、ドット 2 つ $\ddot x$ は 2 階微分 $\dd^2x/\dd t^2$ を表す。
太字（$\bm r$, $\bm v$）はベクトル。

\begin{center}
\begin{tabularx}{\linewidth}{@{}lLl@{}}
\toprule
記号 & 意味 & 単位 \\
\midrule
$t$ & 時間 & \unit{s} \\
$h$（コード上は \code{dt}） & 時間刻み & \unit{s} \\
$g$ & 重力加速度（9.81） & \unit{m/s^2} \\
$L$ & 振り子の長さ & \unit{m} \\
$\theta$ & 鉛直下向きからの角度 & \unit{rad} \\
$\omega=\dot\theta$ & 角速度 & \unit{rad/s} \\
$\omega_0=\sqrt{g/L}$ & 小角近似の角振動数 & \unit{rad/s} \\
$T_0=2\pi/\omega_0$ & 小角近似の周期 & \unit{s} \\
$E$ & 力学的エネルギー（単位質量あたりで書くことが多い） & \unit{J}（または \unit{J/kg}） \\
\midrule
$q_i$ & 一般座標（$i=1,\dots,s$） & 座標による \\
$\dot q_i$ & 一般速度 & \\
$s$ & 自由度の数 & \\
$\Lagr$ & ラグランジアン（ランダウ＝リフシッツの $L$。振り子の長さ $L$ と区別する） & \unit{J} \\
$S$ & 作用 & \unit{J.s} \\
$\delta$ & 変分（経路をわずかに変えたときの変化） & \\
$\bm r,\bm v$ & 位置ベクトル、速度ベクトル & \unit{m}, \unit{m/s} \\
$\bm V$ & 慣性系どうしの相対速度 & \unit{m/s} \\
$T$, $U$ & 運動エネルギー、ポテンシャルエネルギー（周期の $T$ とは文脈で区別） & \unit{J} \\
\midrule
$p_i=\partial\Lagr/\partial\dot q_i$ & 一般化運動量 & 座標による \\
$E$ & エネルギー $\sum p_i\dot q_i-\Lagr$（$T+U$ に一致しない場合がある） & \unit{J} \\
$\bm P$ & 全運動量 & \unit{kg.m/s} \\
$M$, $\bm R$ & 全質量、重心の位置 & \unit{kg}, \unit{m} \\
$\bm L$ & 角運動量（ラグランジアン $\Lagr$ と区別する） & \unit{kg.m^2/s} \\
$k$ & ポテンシャルの同次性の次数 $U(\alpha\bm r)=\alpha^kU(\bm r)$ & \\
\bottomrule
\end{tabularx}
\end{center}

% 見出しと表が別のページに分かれないように、まとめて 1 つの箱にする
\noindent\begin{minipage}{\linewidth}
\paragraph{衝突と散乱（第\ref{ch:coll}章）}\mbox{}\par\smallskip
\begin{center}
\begin{tabularx}{\linewidth}{@{}lLl@{}}
\toprule
記号 & 意味 & 単位 \\
\midrule
$\mu$ & 換算質量 $m_1m_2/(m_1+m_2)$ & \unit{kg} \\
$\varepsilon$ & 崩壊で運動エネルギーになる分 & \unit{J} \\
$\chi$ & 重心系の散乱角 & \unit{rad} \\
$\theta_1,\theta_2$ & 実験室系で入射粒子・標的が出ていく角 & \unit{rad} \\
$\rho$ & 衝突パラメータ（円筒座標の $\rho$ と文脈で区別） & \unit{m} \\
$v_\infty$ & 十分に離れたときの相対速度の大きさ & \unit{m/s} \\
$\sigma$, $\dd\Omega$ & 断面積、立体角 & \unit{m^2}, \unit{sr} \\
$\alpha$ & ポテンシャル $\alpha/r^n$ の係数（正が斥力） & \unit{J.m}（$n=1$） \\
$d=\abs{\alpha}/E$ & 正面衝突の最接近距離 & \unit{m} \\
$\epsilon$ & 静止していた標的に渡るエネルギー & \unit{J} \\
$k$（崩壊の話だけ） & 親の速さと破片の静止系の速さの比 $V/v_0$ & \\
\bottomrule
\end{tabularx}
\end{center}
\end{minipage}
